\section{Preliminary}
In standard reinforcement learning (RL)\cite{sutton1998reinforcement} tasks, the problem is typically modeled as a Markov Decision Process (MDP), denoted by \( \mathcal{M} = \langle \mathcal{S}, \mathcal{A}, \mathcal{T}, R, \gamma, \rho_0 \rangle \). Here, \( \mathcal{S} \) and \( \mathcal{A} \) represent the state and action spaces, respectively, \( \gamma \) is the discount factor, and \( \rho_0 \) is the distribution of initial states. The dynamics of environment are captured by the transition function \( \mathcal{T}: \mathcal{S} \times \mathcal{A} \times \mathcal{S} \to [0, 1] \), which defines the probability of transitioning from one state to another given a specific action. The reward function \( R: \mathcal{S} \times \mathcal{A} \to \mathbb{R} \) quantifies the immediate reward received for taking an action in a particular state.
A trajectory \( \tau = (s_0, a_0, r_0, \dots, s_T, a_T, r_T) \) describes the sequence of states, actions, and rewards encountered by an agent following a decision-making policy \( \pi \). At each time step \( t \), the agent selects an action \( a_t \) according to \( a_t \sim \pi(\cdot | s_t) \), and the environment responds by transitioning to the next state \( s_{t+1} \sim \mathcal{T}(\cdot | s_t, a_t) \).
The state-action value function (Q-function) $Q^\pi(s, a)$ represents the expected cumulative return when taking action $a$ in state $s$ and subsequently following policy $\pi$:
\begin{equation}
    Q^\pi(s, a) = \mathbb{E}_{\tau \sim p(\tau|\pi)} \left[ \sum_{k=0}^{T} \gamma^k R(s_{t+k}, a_{t+k}) \,\bigg|\, s_t = s, a_t = a \right]
\end{equation}
The objective in RL is for the agent to learn a policy \( \pi : \mathcal{S} \to \mathcal{A} \) that maximizes the expected cumulative return defined as:
\begin{equation}
    J(\pi) = \mathbb{E}_{\tau \sim p(\tau|\pi)} \left[ \sum_{t=0}^{T} \gamma^t R(s_t, a_t) \right]
\end{equation}
